\section*{Chapitre \ref{ch:g12:innate-immunity} --- L'immunité innée}

\begin{solution}{exo:g12:innate-immunity:1}
Innée~: présente dès la naissance, agit en quelques heures, reconnaît
des catégories d'intrus à des motifs partagés, répond de façon
identique à chaque fois. Adaptative~: met des jours, reconnaît chaque
intrus individuellement, s'en souvient (et n'existe que chez les
\omterm{def:g10:common-ancestry:bodyplan}{vertébrés}).
\end{solution}

\begin{solution}{exo:g12:innate-immunity:2}
Rougeur et chaleur~: des vaisseaux dilatés qui apportent plus de sang.
Gonflement~: des parois vasculaires perméables qui laissent passer le
plasma dans le tissu. Douleur~: des terminaisons nerveuses
sensibilisées par les \omterm{prop:g12:innate-immunity:inflammation}{médiateurs}.
\end{solution}

\begin{solution}{exo:g12:innate-immunity:3}
Des motifs moléculaires partagés par des catégories de microbes
(fragments de paroi, \omterm{def:g10:chemistry-of-life:families}{acides nucléiques} viraux) et des molécules
libérées par les \omterm{def:g10:cells-common-unit:cell}{cellules} abîmées. Elles libèrent des \omterm{prop:g12:innate-immunity:inflammation}{médiateurs}~:
histamine, prostaglandines, cytokines.
\end{solution}

\begin{solution}{exo:g12:innate-immunity:4}
Contact et liaison par des récepteurs~; englobement par la membrane~;
digestion dans une vésicule au moyen d'\omterm{def:g11:enzymes-and-phenotype:enzyme}{enzymes} et de composés
réactifs~; exposition de fragments protéiques à la surface.
\end{solution}

\begin{solution}{exo:g12:innate-immunity:5}
Ils bloquent l'\omterm{def:g11:enzymes-and-phenotype:enzyme}{enzyme} qui fabrique les prostaglandines, ce qui réduit
le gonflement, la douleur et la fièvre. Ils ne suppriment pas la cause
et ne tuent aucun microbe.
\end{solution}

\begin{solution}{exo:g12:innate-immunity:6}
L'innée vers une demi-journée~; l'adaptative vers le jour 9~; les
courbes se croisent autour des jours 5 et 6.
\end{solution}

\begin{solution}{exo:g12:innate-immunity:7}
Le déclencheur est un motif moléculaire reconnu par les \omterm{def:g10:cells-common-unit:cell}{cellules}
sentinelles, non un intrus vivant. Un sérum physiologique stérile ne
porte ni motif ni lésion~: tout au plus une légère réaction à
l'aiguille.
\end{solution}

\begin{solution}{exo:g12:innate-immunity:8}
Des neutrophiles morts, des bactéries mortes et vivantes, des débris
digérés et du plasma. Les neutrophiles arrivent au fil des heures et
meurent au fil d'une journée~; le pus est ce qui s'accumule une fois
qu'ils ont combattu.
\end{solution}

\begin{solution}{exo:g12:innate-immunity:9}
Le rhume des foins, c'est de l'histamine libérée par les mastocytes
contre du pollen~: bloquer l'histamine supprime la réaction. Dans une
infection, l'histamine ne mène que les premières minutes~; le
gonflement, la douleur et le recrutement ultérieurs reposent sur les
prostaglandines et les cytokines, auxquelles l'antihistaminique ne
touche pas.
\end{solution}

\begin{solution}{exo:g12:innate-immunity:10}
À 12 heures, 100\% contre 55\%~; à 36 heures, 50\% contre 30\%. Le
médicament abaisse le gonflement à tout instant, mais la courbe tombe
toujours à zéro vers 72 heures~: il ne raccourcit pas la réaction.
\end{solution}

\begin{solution}{exo:g12:innate-immunity:11}
Des fibres ligamentaires déchirées et des \omterm{def:g10:cells-common-unit:cell}{cellules} écrasées libèrent
les molécules des \omterm{def:g10:cells-common-unit:cell}{cellules} abîmées, que les récepteurs des \omterm{def:g10:cells-common-unit:cell}{cellules}
sentinelles reconnaissent comme ils le feraient d'un microbe~: la
lésion seule est un déclencheur.
\end{solution}

\begin{solution}{exo:g12:innate-immunity:12}
Sans neutrophiles, les premiers jours sont perdus et des bactéries
ordinaires submergent l'organisme~: la réaction innée est ce qui tient
la ligne d'emblée. Sans \omterm{def:g12:innate-immunity:immunity}{immunité adaptative}, la réaction innée règle
les premières infections, mais celles qui persistent ou se répètent,
et qui exigent des anticorps spécifiques et de la mémoire, finissent
par l'emporter~: la réaction adaptative achève ce que l'innée
contient.
\end{solution}

\begin{solution}{exo:g12:innate-immunity:13}
Un abcès est une inflammation qui combat des bactéries vivantes~:
l'amortir ralentit leur élimination et risque de laisser l'infection
s'étendre. Une arthrite est une inflammation sans microbe, dirigée
contre l'\omterm{def:g10:muscles-and-joints:joint}{articulation} elle-même~: là, la réaction est la maladie, et
la réduire est le traitement.
\end{solution}

\begin{solution}{exo:g12:innate-immunity:14}
Les récepteurs innés sont fixés par les \omterm{def:g10:universal-dna:gene}{gènes} et reconnaissent des
motifs partagés par des catégories entières~; rien, dans une seconde
rencontre, ne les modifie. Le système adaptatif construit des
récepteurs propres à chaque intrus et conserve les \omterm{def:g10:cells-common-unit:cell}{cellules} qui les
ont fabriqués~: il reconnaît l'individu, et il garde cette
reconnaissance.
\end{solution}

\begin{solution}{exo:g12:innate-immunity:15}
L'inflammation est le remède quand elle est dirigée contre un intrus
ou une lésion réels, qu'elle est proportionnée et qu'elle se résout en
quelques jours~: elle nettoie le site et amorce la réparation. Elle
est la maladie quand elle est dirigée contre quelque chose
d'inoffensif (un pollen, les propres \omterm{def:g10:muscles-and-joints:joint}{articulations} du corps),
excessive ou sans fin~: ses \omterm{prop:g12:innate-immunity:inflammation}{médiateurs} et ses \omterm{def:g12:innate-immunity:phagocytosis}{phagocytes} abîment alors
le tissu qu'ils devaient protéger. Ce qui en décide, c'est la cible et
la durée, non le mécanisme, qui est le même dans les deux cas.
\end{solution}

\begin{solution}{pb:g12:innate-immunity:1}
\textbf{1.} Trois doublements en 2 heures~: $\num{10000} \times 8 =
\num{80000}$~; six en 4 heures~: $\num{640000}$.

\textbf{2.} $200 \times 5 = 1000$ par heure --- bien en dessous des
$\num{10000}$ bactéries nouvelles qu'ajoutent les doublements de la
première heure~: la croissance continue.

\textbf{3.} $\num{50000} \times 20 = 10^6$ bactéries.

\textbf{4.} À l'heure 2, il y a moins de $\num{80000}$ bactéries~; les
seuls neutrophiles de la première heure peuvent en retirer un million.
Les \omterm{def:g12:innate-immunity:phagocytosis}{phagocytes} l'emportent dès l'heure 2.

\textbf{5.} En quelques heures, les bactéries ont disparu~; il reste
les neutrophiles morts, les débris et le liquide qui feront le pus, et
des macrophages qui nettoient.

\textbf{6.} Rougeur et chaleur~: l'histamine et les prostaglandines
dilatent les vaisseaux. Gonflement~: l'histamine rend les parois
perméables et le plasma s'infiltre. Douleur~: les prostaglandines
sensibilisent les terminaisons nerveuses.

\textbf{7.} Les vaisseaux perméables sont le chemin par lequel les
neutrophiles et les \omterm{def:g10:chemistry-of-life:families}{protéines} défensives du plasma gagnent le tissu~;
le liquide qui appuie sur les nerfs sensibilisés, c'est la douleur.

\textbf{8.} Chaque battement de cœur pousse plus de sang dans les
vaisseaux dilatés~; le tissu gonflé ne peut pas se dilater davantage,
si bien que chaque pulsation élève la pression sur les nerfs
sensibilisés~: un élancement au rythme du cœur.

\textbf{9.} Des neutrophiles morts, des bactéries mortes et vivantes,
des débris digérés, du plasma.

\textbf{10.} Les motifs ont disparu avec les bactéries et les débris
ont été nettoyés~: plus de motif, plus de reconnaissance par les
sentinelles, plus de \omterm{prop:g12:innate-immunity:inflammation}{médiateurs}~; les vaisseaux reviennent à la
normale.

\textbf{11.} Il réduit la rougeur, la chaleur et la fuite précoces
dues à l'histamine~; non le gonflement et la douleur plus tardifs, dus
aux prostaglandines, ni le recrutement des neutrophiles.

\textbf{12.} Il bloque l'\omterm{def:g11:enzymes-and-phenotype:enzyme}{enzyme} qui fabrique les prostaglandines~: la
douleur, une part du gonflement et la fièvre baissent.

\textbf{13.} La crème éteint la plupart des \omterm{prop:g12:innate-immunity:inflammation}{médiateurs}, si bien que
moins de neutrophiles sont appelés~; avec moins de \omterm{def:g12:innate-immunity:phagocytosis}{phagocytes}, les
bactéries peuvent ne pas être éliminées et l'infection s'étendre.

\textbf{14.} Tous trois agissent sur les signes et aucun ne tue de
bactérie~; le corticoïde, en coupant le recrutement des \omterm{def:g12:innate-immunity:phagocytosis}{phagocytes},
peut faire durer l'infection plus longtemps.

\textbf{15.} Des cytokines parvenues au cerveau, qui élèvent la
consigne de température~; la chaleur en plus accélère les \omterm{def:g12:innate-immunity:phagocytosis}{phagocytes}
et ralentit les bactéries.

\textbf{16.} Les cytokines qu'elles ont captées sur le site, qui
disent aux lymphocytes de quel genre d'intrus il s'agissait~; elles
font leur rapport aux lymphocytes du ganglion lymphatique.

\textbf{17.} Environ une semaine à dix jours. Pour une écharde réglée
en quelques heures par la réaction innée, ils ne seront pas
nécessaires --- mais ils seront fabriqués, et mémorisés.

\textbf{18.} La réaction adaptative est plus rapide et plus forte la
seconde fois, grâce à la mémoire~; la réaction innée est exactement la
même.

\textbf{19.} Les \omterm{def:g11:antibiotic-resistance:antibiotic}{antibiotiques} tuent les bactéries et tiennent la
place des \omterm{def:g12:innate-immunity:phagocytosis}{phagocytes} manquants contre les infections bactériennes~;
ils ne peuvent remplacer ni le nettoyage des débris par les
\omterm{def:g12:innate-immunity:phagocytosis}{phagocytes}, ni leur action contre les champignons, ni le rapport au
ganglion lymphatique.

\textbf{20.} Vers l'heure 2, les \omterm{def:g12:innate-immunity:phagocytosis}{phagocytes} dépassent les bactéries~;
l'histamine et les prostaglandines sont derrière la rougeur, la
chaleur et le gonflement, les prostaglandines derrière la douleur~;
c'est la \omterm{prop:g12:innate-immunity:handover}{cellule dendritique} qui porte le rapport.
\end{solution}
